%! Piecewise Linear Learning Functions — Unit 8, Lesson 8.1 — Group Activity
\documentclass[10pt]{article}
\usepackage{linalg-article}
\usepackage{linalg-boxes}
\newcommand{\TallMath}[1]{$\displaystyle #1\rule[-1.4em]{0pt}{3.2em}$}
\newcommand{\vv}{\mathbf{v}}
\newcommand{\bb}{\mathbf{b}}
\newcommand{\zero}{\mathbf{0}}
\newcommand{\relu}{\mathrm{ReLU}}

\begin{document}

\pageheader{Unit 8, Lesson 8.1}{Group Activity}
\namepartnerperiod

\vspace{0.10in}

\begin{headlinebox}{blush}
\small\textbf{Folding Lines into a Fit.} A \textbf{piecewise-linear} function is a sum of shifted ReLUs:
each $\relu(x-s)$ adds a \emph{fold} at $x=s$, and its coefficient sets how much the slope changes there. On
any piece, the slope is the sum of the weights of the ReLUs already switched on, and $N$ folds give $N+1$
straight pieces. Work with your partner, show each step, and \emph{explain} the shape you get.
\end{headlinebox}

\vspace{0.10in}

\begin{tcolorbox}[colback=white, colframe=black!40, title=\textbf{Tier R --- Remediate},
    fonttitle=\bfseries, arc=1.5mm, left=3mm, right=3mm, top=2mm, bottom=2mm, breakable]
Warm up the two moves: shift a fold, then add ReLUs.
\begin{enumerate}[leftmargin=1.7em, itemsep=8pt]
  \item \textbf{The shift moves the fold.} Evaluate $\relu(x-2)$:
        \[ \relu(0-2)=\blank{0.9cm},\ \ \relu(2-2)=\blank{0.9cm},\ \ \relu(5-2)=\blank{0.9cm}. \]
        This function is flat until $x=\blank{0.8cm}$ (its fold), then rises.
  \item \textbf{A ramp from two ReLUs.} Let $G(x)=\relu(x)-\relu(x-3)$. Fill the table, then answer.
        \[ G(-1)=\blank{0.8cm},\ \ G(0)=\blank{0.8cm},\ \ G(1)=\blank{0.8cm},\ \ G(3)=\blank{0.8cm},\ \ G(5)=\blank{0.8cm}. \]
        Where are the two folds, and what does the graph do after $x=3$? \writeline
\end{enumerate}
\end{tcolorbox}

\begin{tcolorbox}[colback=white, colframe=black!40, title=\textbf{Tier A --- Approaching Proficiency},
    fonttitle=\bfseries, arc=1.5mm, left=3mm, right=3mm, top=2mm, bottom=2mm, breakable]
Now build a tent, place a fold from a bias, and count pieces.
\begin{enumerate}[leftmargin=1.7em, itemsep=9pt]
  \item \textbf{Build a tent.} Let $F(x)=\relu(x)-2\,\relu(x-3)+\relu(x-6)$. Compute $F(3)$, $F(4)$, and
        $F(6)$. Give the slope on each of the four pieces, and name the peak. \writelines{2}
  \item \textbf{The bias sets the fold.} A hidden unit computes $\relu(2x-6)$. Its fold is where
        $2x-6=0$. Solve for that $x$, then evaluate the unit at $x=1$ and $x=5$. \writeline
  \item \textbf{Count the pieces.} A network's hidden layer has $5$ ReLUs. How many folds, and how many
        straight pieces, can its function have? \writeline
\end{enumerate}
\end{tcolorbox}

\begin{tcolorbox}[colback=white, colframe=black!40, title=\textbf{Tier E --- Extension},
    fonttitle=\bfseries, arc=1.5mm, left=3mm, right=3mm, top=2mm, bottom=2mm, breakable]
\textbf{Why a wider layer fits more.} More ReLUs mean more folds --- and more folds mean the graph can
change direction more often.
\begin{enumerate}[leftmargin=1.7em, itemsep=9pt]
  \item Fill in the counting rule:
        \[ \begin{array}{c|c} \text{ReLUs } (N) & \text{straight pieces} \\ \hline
           1 & \blank{0.8cm} \\ 2 & \blank{0.8cm} \\ 3 & \blank{0.8cm} \\ 4 & \blank{0.8cm} \end{array} \]
  \item A data pattern goes \emph{up}, then \emph{down}, then \emph{up} again (two changes of direction).
        What is the fewest number of folds a piecewise-linear function needs to make those two turns? \writeline
  \item In one sentence: why does a \emph{wider} weight matrix $A$ (more rows, so more ReLUs) let a network
        fit more complicated, wigglier data? \writeline
\end{enumerate}
\end{tcolorbox}

\end{document}
